\documentclass{article}

\usepackage{amsmath,amssymb}
\usepackage{xurl}
\usepackage[hidelinks]{hyperref}
\setlength{\emergencystretch}{2em}

% Shared commands for notes in docs/paper-gaps/.

\DeclareUnicodeCharacter{2080}{\ifmmode{}_{0}\else\textsubscript{0}\fi}
\DeclareUnicodeCharacter{2081}{\ifmmode{}_{1}\else\textsubscript{1}\fi}
\DeclareUnicodeCharacter{2082}{\ifmmode{}_{2}\else\textsubscript{2}\fi}
\DeclareUnicodeCharacter{2099}{\ifmmode{}_{n}\else\textsubscript{n}\fi}

\newcommand{\Lean}{\textsc{Lean}}
\ExplSyntaxOn
\NewDocumentCommand{\leanid}{m}{
  \ifmmode
    \textsf{\detokenize{#1}}
  \else
    \group_begin:
    \urlstyle{sf}
    \tl_set:Nn \l_tmpa_tl {#1}
    \tl_replace_all:Nnn \l_tmpa_tl {\_} {_}
    \exp_args:NV \path \l_tmpa_tl
    \group_end:
  \fi
}
\ExplSyntaxOff
\providecommand{\tr}{\operatorname{tr}}
\newcommand{\tnleanIssueBase}{https://github.com/LionSR/TNLean/issues/}
\newcommand{\tnleanPrBase}{https://github.com/LionSR/TNLean/pull/}
\newcommand{\tnleanDocBase}{https://sirui-lu.com/TNLean/docs/}
\newcommand{\ghissue}[1]{\href{\tnleanIssueBase#1}{GitHub issue \##1}}
\newcommand{\ghpr}[1]{\href{\tnleanPrBase#1}{PR \##1}}
\newcommand{\doclink}[1]{\href{\tnleanDocBase#1.html}{\path{#1}}}

% Dirac notation and the identity operator, as in blueprint/src/macros/common.tex.
% \providecommand keeps note-local definitions authoritative.
\usepackage{dsfont}
\providecommand{\ket}[1]{|#1\rangle}
\providecommand{\bra}[1]{\langle #1|}
\providecommand{\braket}[2]{\langle #1 | #2 \rangle}
\providecommand{\ketbra}[2]{|#1\rangle\!\langle #2|}
\providecommand{\Id}{\mathds{1}}
\providecommand{\one}{\mathds{1}}
\providecommand{\C}{\mathbb{C}}
\providecommand{\MN}[1]{M_{#1}(\C)}
\providecommand{\MD}{M_D(\C)}

% External referencing for paper sources.  The build helper writes sanitized
% auxiliary files under build/paper-aux/ and excludes duplicated source labels.
\usepackage{xr}

% Import a generated paper auxiliary file with a caller-supplied label prefix.
% The candidate paths support builds from docs/paper-gaps/, the repository root,
% and build/paper-aux/ itself.
\newcommand{\importpaperaux}[2]{%
  \IfFileExists{../../build/paper-aux/#2.aux}{%
    \externaldocument[#1]{../../build/paper-aux/#2}%
  }{%
    \IfFileExists{build/paper-aux/#2.aux}{%
      \externaldocument[#1]{build/paper-aux/#2}%
    }{%
      \IfFileExists{#2.aux}{%
        \externaldocument[#1]{#2}%
      }{}%
    }%
  }%
}

% General external reference macros
\newcommand{\paperref}[2]{\ref{#1-#2}}
\newcommand{\papereqref}[2]{(\ref{#1-#2})}

% CPSV16 source (1606.00608 / MPDO-22-12-17-2.tex) external document and shortcuts
\importpaperaux{cpsv16-}{1606.00608/MPDO-22-12-17-2-tnlean}
\newcommand{\cpsvref}[1]{\paperref{cpsv16}{#1}}
\newcommand{\cpsveqref}[1]{\papereqref{cpsv16}{#1}}

% Verdict marker: \gapnote{<kind>}{<status>}. Kind is the policy
% classification (clarification, local-correction, scope-restriction,
% unfaithful, false-source, open-gap); status is open, wip (elimination
% underway), resolved, or historical. Machine-read by the site index;
% prints a verdict line.
\newcommand{\gapnote}[2]{\par\noindent\textbf{Verdict:} #1 (#2).\par}


\title{The Even-Dimension Restriction in the Indecomposability\\
of the Breuer--Hall Map (Wolf Example~3.1)}
\author{}
\date{2026-08-06}

\begin{document}
\maketitle

\section{Source statement}

Wolf's Chapter~3, Example~3.1 introduces the Breuer--Hall
map\footnote{\cite[Chapter~3, Example~3.1]{Wolf2012Quantum}; local source:
\path{Notes/WolfNoteTexSource/ch03_positive_not_completely.tex}, lines
344--355. Formalization issue:
\href{https://github.com/LionSR/TNLean/issues/5471}{issue \#5471}.}
\begin{align}
  T_{\mathrm{BH}}(X)=\operatorname{tr}[X]\,\mathbb 1-X-UX^{T}U^{\dagger}.
  \tag{3.19}
\end{align}
Here $U$ may be any antisymmetric contraction satisfying $U^{T}=-U$ and
$U^{\dagger}U\le\mathbb 1$.  The source then states: ``For $d$ even there are
anti-symmetric unitaries
$U^{\dagger}=U^{-1}=-\overline{U}$ \dots\ If $U$ is an anti-symmetric unitary
$T_{\mathrm{BH}}$ is not decomposable.''  The final sentence has no dimension
hypothesis.  Thus the source uses two distinct hypothesis regimes: the map is
defined and proved positive for antisymmetric contractions
$U^{\dagger}U\le\mathbb 1$, whereas its indecomposability assertion concerns
antisymmetric unitaries $U^{\dagger}U=\mathbb 1$.  The formalization proves
indecomposability only in the unitary case; extending that conclusion to general
antisymmetric contractions remains open.

\section{The point at issue}

The blanket statement fails at $d=2$. Every antisymmetric $2\times2$ matrix
is a scalar multiple of
\begin{align}
  J=\begin{pmatrix} 0 & 1 \\ -1 & 0 \end{pmatrix},
\end{align}
so every antisymmetric unitary at $d=2$ is a phase $e^{i\theta}J$, and the
twisted-transpose term is independent of the phase. A direct computation
gives, for $X=(\begin{smallmatrix} a & b \\ c & d\end{smallmatrix})$,
\begin{align}
  JX^{T}J^{\dagger}
  =\begin{pmatrix} d & -b \\ -c & a\end{pmatrix}
  =\operatorname{tr}[X]\,\mathbb 1-X,
\end{align}
so $T_{\mathrm{BH}}$ is the zero map on $\mathcal M_2$. The zero map is
completely positive, hence decomposable. The hypothesis $d>2$ is therefore
necessary, and the source's sentence must be read with the restriction
$d\ge4$ (recall that antisymmetric unitaries exist only for even $d$: taking
determinants in $U^{T}=-U$ gives $\det U=(-1)^{d}\det U$, and $\det U\ne0$).

\section{The corrected statement}

The proved statement is: if $d>2$ and $U$ is an antisymmetric unitary on
$\mathbb C^{d}$, then $T_{\mathrm{BH}}$ is an indecomposable positive
map.\footnote{Formal declarations:
\leanid{Matrix.breuerHallMap\_isIndecomposablePositiveMap} and
\leanid{Matrix.breuerHallMap\_not\_isDecomposablePositiveMap} in
\doclink{TNLean/Channel/BreuerHallIndecomposable}; positivity is
\leanid{Matrix.breuerHallMap\_isPositiveMap} in
\doclink{TNLean/Channel/BreuerHallMap}.}
The source gives no proof; the formalized argument is the standard PPT
detection route. Write $F$ for the swap operator on
$\mathbb C^{d}\otimes\mathbb C^{d}$ and $\ket{\Psi}$ for the $U$-twisted
maximally entangled vector, $\Psi(i,k)=U(i,k)$, with unnormalized projector
$P_{0}=\ket{\Psi}\bra{\Psi}$. The witness state
\begin{align}
  \rho=(\mathbb 1+F)+P_{0}
\end{align}
is a sum of positive semidefinite terms, and antisymmetry with unitarity of
$U$ gives the partial-transpose evaluations
$(\mathbb 1+F)^{T_{1}}=\mathbb 1+P_{1}$ and
$P_{0}^{T_{1}}=(U^{\dagger}\otimes\mathbb 1)(\mathbb 1+F)(U\otimes\mathbb 1)
-\mathbb 1$, where $P_{1}$ is the untwisted maximally entangled projector.
Hence
\begin{align}
  \rho^{T_{1}}=P_{1}+(U^{\dagger}\otimes\mathbb 1)(\mathbb 1+F)(U\otimes\mathbb 1)
  \ge0,
\end{align}
so $\rho$ is a PPT state. Expanding all tensor indices and contracting
against $U^{\dagger}U=\mathbb 1$ yields the carrying identity, valid for
every bipartite $\sigma$,
\begin{align}
  \operatorname{tr}[(T_{\mathrm{BH}}\otimes\mathrm{id})(\sigma)\,P_{0}]
  =\operatorname{tr}(\sigma)-\operatorname{tr}(\sigma P_{0})
  -\operatorname{tr}(F\sigma).
\end{align}
At $\sigma=\rho$ the three traces are $d^{2}+2d$, $d^{2}$, and $d^{2}$ ---
the last two because antisymmetry gives $F\ket{\Psi}=-\ket{\Psi}$ --- so the
pairing equals $d(2-d)<0$ for $d>2$, and
$(T_{\mathrm{BH}}\otimes\mathrm{id})(\rho)$ is not positive semidefinite.
Since decomposable maps preserve positivity on PPT
states,\footnote{Formal declarations:
\leanid{IsDecomposablePositiveMap.tensorMapId\_posSemidef\_of\_isPPT} and
\leanid{not\_isDecomposablePositiveMap\_of\_isPPT\_not\_tensorMapId\_posSemidef}
in \doclink{TNLean/Channel/DecomposablePPT}.}
$T_{\mathrm{BH}}$ is not decomposable.

\section{Conclusion}

The source's indecomposability claim for the Breuer--Hall map is correct
exactly in the dimensions where antisymmetric unitaries exist and $d>2$,
i.e.\ for even $d\ge4$; at $d=2$ the map vanishes. The formalization proves
the corrected statement with the explicit hypothesis $d>2$, and the
blueprint entry records the restriction in its statement.

\bibliographystyle{alpha}
\bibliography{references}

\end{document}